\subsection{Naturality of the scale and operator-valued canonical forms}
\label{rev5:naturalidad-areal}

The generated length makes the compatibility between the positive
realisation, the area operator, and sectorial information explicit.
The marked representation \(\mathcal F(K,\xi)\) retains,
among the relevant data in \(\xi\), the resolution
\(\Omega\), its projectors, the transported return, and the section
\(L_*\). The generated regional coordinates, the action, and the constitutive
channels retain their preceding genealogy. The marked inverse
therefore recovers every argument of
\eqref{rev5:longitud-generada} and \eqref{rev5:G-area}.
In particular, the previously constructed geometric reader of
\(\alpha^{\mathrm{geom}}\) determines the expressions
\[
 L_\alpha^{\mathrm{geom}}
 =L_*(\alpha^{\mathrm{geom}})^{16}\frac{5759}{23040},
 \qquad
 G^{\mathrm{geom}}
 =\frac{(c^{\mathrm{geom}})^3
 (L_\alpha^{\mathrm{geom}})^2}{\hbar_{\rm ret}^{\mathrm{geom}}}.
\]
The speed and action readers are obtained by the same composition
with the marked inverse. These expressions retain the return data
and the radial reference section.

\begin{corollary}[Preservation under changes of chart and inherited refinement]
\label{rev5:conservacion}
Positive changes of chart preserving the marked point, and
refinements recovering the initial record, preserve
\(L_\alpha\), \(G\), \(\gamma_{\rm HMT}\), and
\(\widehat{\mathcal A}=L_\alpha^2a_\partial N\), provided
they also transport the declared sectorial and dimensional data.
The normalised spectrum
\(\widehat{\mathcal A}/L_\alpha^2=a_\partial N\), its
multiplicities, and the state \(\rho_A(\Delta)\) remain unchanged.
\end{corollary}
\begin{proof}
Theorem~\ref{thm:compat-algebra-lectores} transports each reader
by composition with the inverse. Mutations of the same point
reconstruct the same record. Under inherited refinement, summing
the subdivided separations reconstructs each initial coordinate.
Tensor extension of the inscription preserves \(r_\Omega\),
while its remaining arguments are fixed by the markings.
Substitution therefore preserves the length, the gravitational
coefficient, and the area operator. The occupation projectors
preserve the multiplicity polynomial \((1+z+z^2)^{36}\) and
the normalised expression of the state.
\end{proof}

The inherited refinement considered here retains the same space
of thirty-six boundary links. An extension introducing a different
physical boundary space requires specified intertwining maps for
\(N_{90}\) and \(N_{120}\), together with the transport of the
declared sectorial data.

Let \(P\) now be a rank-one polytope constructed in this manuscript,
with an oriented subdivision \(P=\bigcup_\sigma P_\sigma\)
having the same support, multiplicity one, and compatible faces.
The same sectorial fibre \(\mathscr H_A\) is fixed over all
its cells. The operator-valued top-degree form with the
dimension of area is defined by
\[
 \boldsymbol\Omega^{\mathcal A}_P
 =\widehat{\mathcal A}\otimes\Omega_P.
\]
Here \(\widehat{\mathcal A}\) is constant with respect to the
differential coordinates of the polytope: it is determined by the
fixed HMT state, whereas \(\Omega_P\) describes its geometric support.
The coefficient \(\widehat{\mathcal A}\) is positive semidefinite.
Its component in the vacuum sector \(N=0\) vanishes. On an
eigenspace with occupation \(n>0\), division by
\(nL_\alpha^2a_\partial\) recovers the canonical form
\(\Omega_P\) multiplied by the identity on that eigenspace.

\begin{proposition}[Areal additivity and boundary residues]
\label{rev5:forma-areal}
In the common sectorial trivialisation,
\begin{equation}
 \boldsymbol\Omega^{\mathcal A}_P
 =\sum_\sigma\widehat{\mathcal A}\otimes\Omega_{P_\sigma},
 \qquad
 \operatorname{Res}_F\boldsymbol\Omega^{\mathcal A}_P
 =\widehat{\mathcal A}\otimes\Omega_F.
\end{equation}
These identities persist under successive subdivisions and
iterated residues. For a fixed state \(\rho\), taking the trace
gives \(\Tr(\rho\widehat{\mathcal A})\Omega_P\) and preserves
the same scalar identities.
\end{proposition}
\begin{proof}
The operator space of \(\mathscr H_A\) is finite-dimensional.
In any basis, the assertion is the additivity identity for
canonical forms multiplied by each constant matrix entry of
\(\widehat{\mathcal A}\). Both residue and trace are linear.
Cancellation at interior facets retains their orientations and
the same operator entry on both sides. Iteration applies these
operations to each subdivision and each flag of faces.
\end{proof}

If distinct cells carry regular operator-valued coefficients
\(A_\sigma\), the defect of a pair along its common face is
\((A_\sigma|_F-A_\tau|_F)\otimes\Omega_F\).
Equality of the operator-valued restrictions to the boundary
provides the gluing condition, under the simple-pole hypotheses
of Theorem~\ref{thm:compat-pegado}. The global sum must also retain
every contribution along each ambient divisor and the control
of exterior poles. Dimensional normalisation, residue cancellation,
and coverage of the support therefore have distinct and
compatible criteria.

The finite modular identity likewise admits a direct expression
in physical areas \(\mathcal A_m=L_\alpha^2a_\partial N_m\):
\begin{equation}
 s(\sigma)-s(\rho_0)
 =\sum_{m=90,120}\frac{m\Delta_0}{L_\alpha^2a_\partial}
   \Tr[(\sigma-\rho_0)\mathcal A_m]
  -D(\sigma\Vert\rho_0).
 \label{rev5:modular-fisica}
\end{equation}
This is \eqref{rev4:ae:ley-finita-formula} after substitution
of the constructed length. The angular factor \(a_0\), contained
in \(a_\partial=\gamma_{\rm HMT}a_0\), is retained in full.
When dimensional sections satisfying \(L'_\alpha=sL_\alpha\)
are compared while preserving the same state and the same
dimensionless data, areas are multiplied by \(s^2\) and their
modular coefficients by \(s^{-2}\). Entropy and relative entropy
remain unchanged. This comparison of sections does not constitute
a second solution for the same fixed longitudinal data.

For a spherical horizon of radius \(R_h=L_\alpha\zeta_h\),
the geometric area is \(4\pi_{\rm HMT}L_\alpha^2\zeta_h^2\).
Its identification with the expectation of the sectorial operator,
when required for the same cross-section, imposes the explicit
condition
\[
 a_\partial\Tr(\rho N)=4\pi_{\rm HMT}\zeta_h^2.
\]
The horizon geometry and the boundary state must satisfy this
relation through their respective constructions. The factor
\(\zeta_h\) characterises that horizon; the longitudinal scale
of the coupling remains \(L_\alpha\).
